\expandafter\ifx\csname CAPETINCLUDEMODE\endcsname\relax
\documentclass[11pt,a4paper]{article}
% preamble.tex — packages et commandes communes a toutes les fiches CAPET SII IE
\input{glyphtounicode}
\pdfgentounicode=1
\usepackage{cmap}
\usepackage[T1]{fontenc}
\usepackage[utf8]{inputenc}
\usepackage[french]{babel}
\usepackage[a4paper, top=1.6cm, bottom=1.6cm, left=1.4cm, right=1.4cm]{geometry}
\setlength{\headheight}{14.5pt}
\usepackage{amsmath, amssymb, amsfonts}
\usepackage{siunitx}
\sisetup{locale=FR, output-decimal-marker={,}, inter-unit-product=\ensuremath{{}\cdot{}}}
\usepackage{array, booktabs, multirow, tabularx}
\usepackage{xcolor}
\usepackage[most, breakable]{tcolorbox}
\usepackage{multicol}
\usepackage{enumitem}
\usepackage{mathtools}
\usepackage{lmodern}
\usepackage{microtype}
\usepackage{fancyhdr}
\usepackage{tikz}
\usetikzlibrary{positioning,arrows.meta}
\usepackage{lastpage}

% ── Couleurs par domaine ──────────────────────────────────────────────
\definecolor{colElec}{RGB}{26, 82, 148}       % Électrotechnique — bleu foncé
\definecolor{colElecLight}{RGB}{210, 225, 245} % fond bleu clair
\definecolor{colAuto}{RGB}{34, 120, 66}        % Automatique — vert
\definecolor{colMeca}{RGB}{185, 90, 20}        % Mécanique — orange
\definecolor{colTherm}{RGB}{160, 40, 40}       % Thermique — rouge
\definecolor{colQT}{RGB}{232, 244, 233}        % fond question type (vert très léger)
\definecolor{colJury}{RGB}{255, 243, 205}      % fond jury (jaune ambre)
\definecolor{colJuryFrame}{RGB}{200, 130, 0}   % bordure jury
\definecolor{colGrayLight}{RGB}{245, 245, 247} % fond gris neutre
\definecolor{colSeparator}{RGB}{180, 180, 180} % gris séparateur

% ── Style tcolorbox de section principale ────────────────────────────
\tcbset{
  ficheSection/.style={
    breakable,
    enhanced,
    colback=#1!7!white,
    colframe=#1,
    colbacktitle=#1,
    coltitle=white,
    fonttitle=\bfseries\normalsize,
    left=4pt, right=4pt, top=3pt, bottom=4pt,
    before skip=6pt, after skip=6pt,
  }
}

% ── Style encadré question type ──────────────────────────────────────
\tcbset{
  questiontype/.style={
    breakable,
    enhanced,
    colback=colQT,
    colframe=colElec!55!black,
    colbacktitle=colElec!55!black,
    coltitle=white,
    fonttitle=\bfseries\small,
    left=4pt, right=4pt, top=3pt, bottom=3pt,
    before skip=5pt, after skip=5pt,
  }
}

% ── Style encadré jury ───────────────────────────────────────────────
\tcbset{
  jury/.style={
    breakable,
    enhanced,
    colback=colJury,
    colframe=colJuryFrame,
    fonttitle=\bfseries\normalsize,
    left=4pt, right=4pt, top=3pt, bottom=3pt,
  }
}

% ── Mise en page ─────────────────────────────────────────────────────
\pagestyle{fancy}
\fancyhf{}
\renewcommand{\headrulewidth}{0.4pt}
\renewcommand{\footrulewidth}{0.3pt}

\setlist[itemize]{leftmargin=1.2em, itemsep=1pt, topsep=2pt, parsep=0pt}
\setlist[enumerate]{leftmargin=1.2em, itemsep=3pt, topsep=2pt, parsep=0pt}

\setlength{\parskip}{2pt}
\setlength{\columnsep}{10pt}

\def\CAPETBEGINFILE{\begin{document}}
\def\CAPETENDFILE{\end{document}}
\else
\def\CAPETBEGINFILE{}
\def\CAPETENDFILE{}
\fi

\fancyhead[L]{\textbf{\textcolor{colElec}{T15}} - Stockage d'energie electrique}
\fancyhead[R]{\small\textit{CAPET SII IE - Session 2026}}
\fancyfoot[C]{\small\thepage\,/\,\pageref{LastPage}}

\CAPETBEGINFILE

% -----------------------------------------------------------------------
%  BLOC EN-TeTE
% -----------------------------------------------------------------------

\begin{tcolorbox}[
  enhanced, breakable,
  colback=colElec,
  colframe=colElec,
  left=6pt, right=6pt, top=5pt, bottom=5pt,
  before skip=0pt, after skip=4pt
]
  {\color{white}\bfseries\large T15 \quad Stockage d'energie electrique}\\[2pt]
  {\color{colElec!30!white}\small
    \textbf{Domaine :} electrotechnique - Systemes energetiques \quad
    \textbf{Presence :} \textbf{7/13} sujets (2013-2026) - theme \emph{frequent}
  }\\[1pt]
  {\color{colElec!20!white}\footnotesize
    Sous-themes : capacite batterie (Aeh) - association serie/parallele de cellules - taux de decharge (DoD) - autonomie - STEP micro-hydraulique - technologies Li-Ion / NMC / LFP - BMS et SOC
  }
\end{tcolorbox}

% -----------------------------------------------------------------------
%  1. RAPPELS DE COURS
% -----------------------------------------------------------------------

\begin{tcolorbox}[
  ficheSection=colElec,
  title={1. Rappels de cours},
]

% -- 1.1 Capacite et energie d'une batterie ---------------------------
\textbf{1.1 Capacite et energie d'une batterie}

\smallskip
\[
  E = U_{nom} \times C_b
  \qquad;\qquad
  C_b = \frac{E_{necessaire}}{U_{nom}}
  \qquad;\qquad
  E_{utile} = E_{totale} \times (1 - DoD_{min})
\]
$C_b$ en Aeh, $E$ en Weh, $U_{nom}$ en V. Le taux de decharge minimum $DoD_{min}$ protege les cellules (ex : $DoD_{min} = 20\,\%$ - energie utile = 80\,\% de l'energie totale).

\medskip\hrule\vspace{4pt}

% -- 1.2 Association de cellules ---------------------------------------
\textbf{1.2 Association serie / parallele de cellules}

\smallskip
\begin{center}
\renewcommand{\arraystretch}{1.4}
\begin{tabular}{llll}
\toprule
\textbf{Association} & \textbf{Tension resultante} & \textbf{Capacite resultante} & \textbf{Formule} \\
\midrule
Serie & $N_s \times U_{cell}$ & $C_{cell}$ inchangee & $N_s = U_{bat}/U_{cell}$ \\
Parallele & $U_{cell}$ inchangee & $N_p \times C_{cell}$ & $N_p = C_{bat}/C_{cell}$ \\
\bottomrule
\end{tabular}
\end{center}
$N_{total} = N_s \times N_p$.

\medskip\hrule\vspace{4pt}

% -- 1.3 Autonomie -----------------------------------------------------
\textbf{1.3 Autonomie et calcul de capacite}

\smallskip
\[
  t_{autonomie} = \frac{E_{utile}}{P_{consommee}}
  \quad\text{(h)}
  \qquad;\qquad
  d_{autonomie} = t_{autonomie} \times v
  \quad\text{(km)}
\]

\medskip\hrule\vspace{4pt}

% -- 1.4 STEP - stockage hydraulique ----------------------------------
\textbf{1.4 STEP micro-hydraulique (Station de Transfert d'energie par Pompage)}

\smallskip
\[
  E_p = m\,g\,h = \rho_{eau}\,V_{utile}\,g\,h
  \qquad;\qquad
  E_{elec} = E_p \times \eta_{turbinage}
  \qquad;\qquad
  \eta_{st} = \frac{E_{elec\,turbinage}}{E_{elec\,pompage}}
\]
$\rho_{eau} = 1000\,\mathrm{kg\!\cdot\!m^{-3}}$, $g = 9{,}81\,\mathrm{m\!\cdot\!s^{-2}}$.
Duree de remplissage : $T_{remplissage} = V_{utile}/Q_{pompe}$.

\medskip\hrule\vspace{4pt}

% -- 1.5 Technologies de cellules -------------------------------------
\textbf{1.5 Technologies de cellules}

\smallskip
\begin{center}
\renewcommand{\arraystretch}{1.2}
\begin{tabular}{llll}
\toprule
\textbf{Technologie} & \textbf{$U_{nom}$} & \textbf{energie massique} & \textbf{Application typique} \\
\midrule
Li-Ion / NMC & 3{,}5-3{,}7\,V & 150-250 Wehekg$^{-1}$ & Bus, VE, robot \\
Li-FePO4 (LFP) & 3{,}2\,V & 100-160 Wehekg$^{-1}$ & Stockage stationnaire \\
Plomb-acide & 2{,}0\,V & 30-50 Wehekg$^{-1}$ & Demarrage, bateau \\
Supercondensateur & 2{,}5-3{,}0\,V & 5-20 Wehekg$^{-1}$ & Tampon puissance \\
\bottomrule
\end{tabular}
\end{center}

\end{tcolorbox}

% -----------------------------------------------------------------------
%  2. FORMULAIRE SYNTHeTIQUE
% -----------------------------------------------------------------------

\begin{tcolorbox}[
  ficheSection=colElec,
  title={2. Formulaire synthetique},
]
\small
\begin{multicols}{2}
\renewcommand{\arraystretch}{1.45}
\begin{tabular}{@{}p{4.5cm}ll@{}}
\toprule
\textbf{Formule} & \textbf{Usage} & \textbf{U.} \\
\midrule
\multicolumn{3}{@{}l}{\footnotesize\textit{Batterie}} \\
$E = U_{nom} \cdot C_b$ & energie & Weh \\
$C_b = E/U_{nom}$ & Capacite & Aeh \\
$E_{utile} = E(1-DoD_{min})$ & Utile & Weh \\
$t = E_{utile}/P$ & Autonomie & h \\
$N_s = U_{bat}/U_{cell}$ & Serie & - \\
$N_p = C_{bat}/C_{cell}$ & Parallele & - \\
\bottomrule
\end{tabular}

\columnbreak

\renewcommand{\arraystretch}{1.45}
\begin{tabular}{@{}p{4.5cm}ll@{}}
\toprule
\textbf{Formule} & \textbf{Usage} & \textbf{U.} \\
\midrule
\multicolumn{3}{@{}l}{\footnotesize\textit{STEP}} \\
$E_p = \rho V g h$ & energie pot. & J \\
$E_{elec} = E_p \cdot \eta_{turb}$ & Turbinage & Weh \\
$\eta_{st} = E_{turb}/E_{pomp}$ & Rendement & - \\
$T = V_{utile}/Q$ & Remplissage & h \\
\midrule
\multicolumn{3}{@{}l}{\footnotesize\textit{Donnees typiques}} \\
$U_{NMC} = 3{,}5\,\mathrm{V}$ & Urbanloop & - \\
$\eta_{STEP} \approx 58\,\%$ & Micro-STEP & - \\
\bottomrule
\end{tabular}
\end{multicols}
\normalsize
\end{tcolorbox}

% -----------------------------------------------------------------------
%  3. QUESTIONS TYPES
% -----------------------------------------------------------------------

\begin{tcolorbox}[
  ficheSection=colElec,
  title={3. Questions types - extraites des sujets reels},
]

% -- QT1 -------------------------------------------------------------

% Mini-procedure de resolution type copie concours
\begin{tcolorbox}[questiontype,title={Procedure type copie concours (5 etapes)}]
\begin{enumerate}[leftmargin=1.5em]
  \item \textbf{Qualifier le modele} : hypotheses, domaine de validite, grandeurs d'entree/sortie.
  \item \textbf{Ecrire les lois} : equations de base avec symboles et unites SI.
  \item \textbf{Isoler l'inconnue} : expression litterale, puis application numerique tracee.
  \item \textbf{Controler l'ordre de grandeur} : coherence physique et unite finale.
  \item \textbf{Conclure concours} : phrase courte orientee decision technique.
\end{enumerate}
\end{tcolorbox}

% Visuel de methode : chemin de resolution CAPET
\begin{center}
\begin{tikzpicture}[font=\small, node distance=7mm and 6mm, >=Stealth]
  \node[draw, rounded corners, fill=colElec!8, inner sep=4pt] (n1) {Notions};
  \node[draw, rounded corners, fill=colElec!8, right=of n1, inner sep=4pt] (n2) {Methode};
  \node[draw, rounded corners, fill=colElec!8, right=of n2, inner sep=4pt] (n3) {Application annale};
  \node[draw, rounded corners, fill=colElec!8, right=of n3, inner sep=4pt] (n4) {Verification jury};
  \draw[->, thick, colElec] (n1) -- (n2);
  \draw[->, thick, colElec] (n2) -- (n3);
  \draw[->, thick, colElec] (n3) -- (n4);
\end{tikzpicture}
\end{center}
\begin{tcolorbox}[questiontype,
  title={QT1 - Association serie/parallele et autonomie - bus Hycity \hfill\normalfont\footnotesize E2 appl. 2026 (Q A3-A5)}
]
\textbf{enonce.}
Batterie bus Hycity : $U_{bat} = 600\,\mathrm{V}$, $E_{bat} = 128{,}4\,\mathrm{kW\!\cdot\!h}$, cellules NMC HpCO ($U_{cell} = 3{,}5\,\mathrm{V}$, $C_{cell} = 53{,}5\,\mathrm{A\!\cdot\!h}$), decharge min 20\,\%. Calculer $N_s$ (cellules serie par module), $C_{totale}$, puis l'autonomie - $v = 40\,\mathrm{km\!\cdot\!h^{-1}}$ pour $P_{moy} = 22\,\mathrm{kW}$.

\tcbline

\textbf{Corrige.}
\[
  N_s = \frac{600/2}{3{,}5} = \frac{300}{3{,}5} \approx 86\,\text{cellules/module} \quad (2\,\text{modules en serie} \times 2\,\text{branches})
\]
\[
  C_{totale} = \frac{E_{bat}}{U_{bat}} = \frac{128{,}4 \times 10^3\,\mathrm{W\!\cdot\!h}}{600\,\mathrm{V}} = 214\,\mathrm{A\!\cdot\!h}
\]
\[
  E_{utile} = 128{,}4 \times (1 - 0{,}20) = 102{,}72\,\mathrm{kW\!\cdot\!h}
  \quad;\quad
  t = \frac{102{,}72}{22} = 4{,}67\,\mathrm{h}
  \quad;\quad
  d = 4{,}67 \times 40 = \mathbf{186{,}8\,\mathrm{km}}
\]
\textbf{Conclusion :} le bus peut rentrer au depet depuis n'importe quel point du trajet type. \checkmark
\end{tcolorbox}

\vspace{4pt}

% -- QT2 -------------------------------------------------------------

\begin{tcolorbox}[questiontype,
  title={QT2 - Dimensionnement batterie - taux de decharge et nombre de cellules \hfill\normalfont\footnotesize Rapport jury 2020 (Q20-Q22)}
]
\textbf{enonce.}
Un robot effectue 160 trajets de $0{,}6\,\mathrm{W\!\cdot\!h}$ chacun. Taux de decharge max : 70\,\%. Tension nominale batterie : 48\,V. Cellules Li-Ion : $U_{cell} = 3{,}7\,\mathrm{V}$, $C_{cell} = 5{,}2\,\mathrm{A\!\cdot\!h}$, masse 92\,g. Calculer $C_b$, puis la structure et verifier la masse (limite 5\,kg).

\tcbline

\textbf{Corrige.}
\[
  E_{necessaire} = \frac{160 \times 0{,}6}{0{,}7} = 137{,}1\,\mathrm{W\!\cdot\!h}
  \qquad;\qquad
  C_b = \frac{137{,}1}{48} = 2{,}86\,\mathrm{A\!\cdot\!h}
\]
\[
  N_s = \frac{48}{3{,}7} = 12{,}97 \to 13 \quad;\quad N_p = \left\lceil\frac{2{,}86}{5{,}2}\right\rceil = 1
  \quad (C_{reelle} = 5{,}2\,\mathrm{A\!\cdot\!h} > 2{,}86\,\mathrm{A\!\cdot\!h})
\]
En pratique le sujet donne $N_p = 4$ branches pour avoir la capacite voulue.
$N_{total} = 13 \times 4 = 52\,\text{cellules}$. Masse = $52 \times 0{,}092 = 4{,}784\,\mathrm{kg} < 5\,\mathrm{kg}$. \checkmark
\end{tcolorbox}

\vspace{4pt}

% -- QT3 -------------------------------------------------------------

\begin{tcolorbox}[questiontype,
  title={QT3 - energie potentielle et rendement STEP \hfill\normalfont\footnotesize E2 disc. 2021 (Q5-Q10), Rapport jury 2021}
]
\textbf{enonce.}
Reservoir : volume total 1000\,m$^3$, volume utile 90\,\%. Hauteur de chute $h = 130\,\mathrm{m}$. Rendement turbinage : 78\,\%. Puissance pompe : 76\,kW, duree remplissage : 5,625\,h. Calculer $E_p$, $E_{elec\,turbinage}$ et le rendement global $\eta_{st}$.

\tcbline

\textbf{Corrige.}
\[
  V_{utile} = 0{,}90 \times 1000 = 900\,\mathrm{m^3}
  \quad;\quad
  m = 1000 \times 900 = 9 \times 10^5\,\mathrm{kg}
\]
\[
  E_p = m\,g\,h = 9 \times 10^5 \times 9{,}81 \times 130 = 1{,}147 \times 10^9\,\mathrm{J} = \frac{1{,}147 \times 10^9}{3{,}6 \times 10^6} = \mathbf{318{,}8\,\mathrm{kW\!\cdot\!h}}
\]
\[
  E_{elec\,turbinage} = 318{,}8 \times 0{,}78 = \mathbf{248{,}7\,\mathrm{kW\!\cdot\!h}}
\]
\[
  E_{pompage} = 76 \times 5{,}625 = 427{,}5\,\mathrm{kW\!\cdot\!h}
  \quad;\quad
  \eta_{st} = \frac{248{,}7}{427{,}5} = \mathbf{58{,}2\,\%}
\]
\textbf{Conclusion :} rendement faible mais technologie ecologique ; complementaire aux panneaux PV (production pendant les heures creuses stockee pour la nuit).
\end{tcolorbox}

\end{tcolorbox}

% -----------------------------------------------------------------------
%  4. ATTENTES DU JURY
% -----------------------------------------------------------------------

\begin{tcolorbox}[
  jury,
  title={4. Attentes du jury - Erreurs recurrentes et points de vigilance},
]

\begin{enumerate}[label=\textbf{\arabic*.}, leftmargin=*, itemsep=5pt]

  \item \textbf{Confusion Weh et Aeh.}
  La capacite s'exprime en Aeh, l'energie en Weh. La conversion est $E = U_{nom} \times C_b$. Ne pas additionner des Aeh et des Weh.

  \item \textbf{Oubli du taux de decharge.}
  Si le constructeur impose $DoD_{max} = 70\,\%$, l'energie utile n'est que 70\,\% de l'energie totale. Rapport jury 2020 : "une critique quantitative etait attendue afin de verifier que l'installation ne depasse pas 70\,\% de decharge".

  \item \textbf{Erreur sur l'association serie/parallele.}
  En \emph{serie} : les tensions s'additionnent, la capacite reste celle d'une cellule. En \emph{parallele} : les capacites s'additionnent, la tension reste celle d'une cellule. Ne jamais multiplier les deux simultanement.

  \item \textbf{Densite de l'eau dans la formule STEP.}
  Utiliser $E_p = \rho V g h$ avec $\rho_{eau} = 1000\,\mathrm{kg\!\cdot\!m^{-3}}$. Oublier la densite donne un resultat 1000 fois trop faible.

  \item \textbf{Rendement STEP = cycle complet.}
  Calculer $\eta_{st} = E_{turbinage}/E_{pompage}$ sur le cycle complet, pas uniquement sur la phase de turbinage (qui donnerait 100\,\%). Citer la valeur numerique (~58\,\%) ET la conclusion contextuelle (ecologique, pas tres efficace).

  \item \textbf{Reponses qualitatives sans chiffres.}
  Rapport jury 2020 : "des reponses trop evasives sans lien avec les documents techniques". Toujours appuyer la conclusion (autonomie suffisante ) par un calcul comparant $E_{disponible}$ et $E_{necessaire}$.

\end{enumerate}

\medskip
\begin{tcolorbox}[
  colback=colJuryFrame!12!colJury,
  colframe=colJuryFrame,
  left=4pt, right=4pt, top=2pt, bottom=2pt,
  fontupper=\small\itshape
]
Rapport jury 2020 : "Le jury deplore que des problematiques simples, comme le dimensionnement d'une batterie, ne soient pas maetrisees. --\\
Rapport jury 2021 : "La conclusion quant aux performances et - l'interet du procede de stockage par pompage devait mentionner le faible rendement (58\,\%) ET l'avantage environnemental par rapport au stockage par batteries."
\end{tcolorbox}

\end{tcolorbox}

% -----------------------------------------------------------------------
%  PIED - PReSENCE DANS LES SUJETS
% -----------------------------------------------------------------------

\begin{tcolorbox}[
  enhanced, breakable,
  colback=colGrayLight,
  colframe=colSeparator,
  left=4pt, right=4pt, top=3pt, bottom=3pt,
  before skip=4pt, after skip=0pt,
  fonttitle=\small\bfseries,
  title={Presence dans les sujets analyses (2013-2026)},
]
\begin{center}
\footnotesize\renewcommand{\arraystretch}{1.15}
\begin{tabular}{lccccccccccccc}
\toprule
\textbf{Annee} & 13 & 14 & 15 & 16 & 17 & 18 & 19 & 20 & 21 & 22 & 23 & 24 & 25 \\
\midrule
E1 (commune)   & - & - & - & - & - & - & - & - & - & - & - & - & - \\
E2 discipl. IE & - & - & - & - & - & - & - & $\checkmark$ & $\checkmark$ & - & $\checkmark$ & $\checkmark$ & $\checkmark$ \\
E2 appl. comm. & - & - & - & - & - & - & - & - & - & - & - & - & - \\
BAC3 / E2 2026 & - & - & - & - & - & - & - & - & - & - & - & - & - \\
\bottomrule
\end{tabular}
\end{center}
\smallskip
\footnotesize\textbf{Note :} la session 2026 inclut un sujet BAC3 (bus electrique Hycity) qui traite le stockage Li-Ion 600\,V / 128{,}4\,kWeh en detail (Q A3-A5).\\
\textbf{Sous-themes frequents :}
capacite batterie en Aeh (5/13) - association serie/parallele (4/13) - taux de decharge DoD (4/13) - energie potentielle STEP (2/13) - rendement cycle complet (2/13)

\medskip
\textbf{Ordres de grandeur (reperes rapides).}
\begin{itemize}[leftmargin=1.3em]
  \item Reseau BT FR : tension phase-neutre typique \SI{230}{\volt}, tension phase-phase typique \SI{400}{\volt}, frequence \SI{50}{\hertz}.
  \item Courants d'etude frequents en sujet : de quelques \si{\ampere} a quelques dizaines de \si{\ampere}.
  \item En conversion d'energie, un rendement attendu reste generalement entre 0{,}8 et 0{,}98 selon la technologie.
\end{itemize}

\medskip
\textbf{Synthese express (a refaire sans support).}
\begin{enumerate}[leftmargin=1.4em]
  \item Poser les hypotheses du modele et les conventions de signe.
  \item Ecrire 2 a 4 relations cles avec unites.
  \item Realiser un calcul numerique complet et verifier l'ordre de grandeur.
  \item Justifier le resultat par un argument technique de niveau CAPET.
\end{enumerate}

\end{tcolorbox}

\CAPETENDFILE









